\documentclass{article}

\usepackage[english]{babel}

\usepackage[
    letterpaper,
    top=2cm,
    bottom=2cm,
    left=3cm,
    right=3cm,
    marginparwidth=1.75cm
]{geometry}

\usepackage{amsmath}
\usepackage{graphicx}
\usepackage[colorlinks=true, allcolors=blue]{hyperref}
\usepackage{amsthm}
\usepackage{csquotes}
\usepackage{algorithm, algpseudocode, float}
\usepackage{enumitem}
\usepackage{setspace}

\theoremstyle{definition}
\newtheorem{definition}{Definition}[section]

\title{COSC 473 Reading Responses}
\author{
    Max Moir \\
    \textit{University of Canterbury} \\ 
    Christchurch, New Zealand \\ 
    \texttt{maxwmoir@gmail.com}
}

\begin{document}

\maketitle

\section{Introduction}

These are the reading responses I wasn't able to turn in due to my clash with the start of the COSC473 lectures.
I believe this covers every week that I didn't submit a response in class.

\subsection{Week 4}

\textit{IPFS was designed to be a decentralized, content-addressed, peer-to-peer foundation for the future of the web.}
\textit{Based on the readings, discuss the following:}
\textit{To what extent does the current reality of IPFS live up to its original idea?}
\textit{Consider one way IPFS fulfills its vision and one way it falls short.}
\\
\\
\noindent
One way that IPFS fulfills its vision is by providing a practical distributed data storage 
application.
This is self evident, as the paper details widespread usage in terms of individual retrievals and 
integration in third party services.
The paper shows presence in 2715 Autonomous Systems and 152 countries, 
which quantifies the fulfillment of this goal for adoption.
Additionally, this ensures IPFS is resistant to single points of failure.
On the other hand, IPFS struggles with a high churn rate, and lower performance when contrasted with 
its centralised peers. The paper states that only 2.5\% of peers stay online for more than 24 hours. 
Additionally, the top 10 ASes contain almost 65\% of all peers.
These are likely symptoms of the lack of reliance on centralised servers, 
but show that IPFS is more consolidated than it intended to be.

\subsection{Week 5}

\textit{In a centralised system, responsibility for knowing where data is concentrated, but Kademlia distributes it.}
\textit{Kademlia is designed so that no one can control it. When is that a strength, and when might it be a problem?}
\textit{Think about where else in life or society we rely on local knowledge and word-of-mouth to make things work without a central coordinator.}
\\
\\
\noindent
The decentralized nature of Kademlia is a strength because it doesn't require a central coordinator.
This means that individual peers communicate with each other to find the location of data, 
and this cannot be limited by a governing body.
A lack of central control means that applications like Kademlia are invulnerable to censorship,
and cannot be limited by a malicious authority.
An example similar to this in real life may be distributing information in opposition to an oppressive government.
Here word of mouth is better than a central coordinator, because a coordinator might become an easy target.
Additionally, if the central coordinator were to be shut down, it would render the entire service useless. 
On the other hand, this lack of authority leads to a trustless environment.
This means moderation is difficult, if possible, and this may be taken advantage of by unsavoury parties.

\subsection{Week 6}

\textit{What kinds of trust are present in Bitcoin, and how do they compare to the forms of trust in earlier systems or traditional financial institutions?}
\textit{How does Bitcoin establish agreement on a shared history of transactions, and how does that compare to Lamport's way of reasoning about order in distributed systems?}
\\
\\
\noindent
The trust required for earlier systems and financial institutions centred primarily around the central authority.
For example, in order to interact with traditional financial institutions, a user would need to trust the central server of, for example, a bank.
This differs in Bitcoin, where events are recorded by a collection of untrusted nodes.
The trust in this decentralized example is moved to the cryptographic algorithms and distributed consensus.
The user must be sure that there are more honest nodes overall in the system than coordinated attackers.
Moreover, the underlying protocols must be trusted to be secure, otherwise the funds in the system would be vulnerable.
There must also be the consensus that tokens are valuable, or else the system would not be worthwhile.
This differs from Lamport's way of reasoning about distributed systems, because it doesn't rely on logical ordering through clocks.
Bitcoin as a distributed system relies on consensus between nodes, with all transactions publicly available and their validity determined by probabilistic consensus.

\subsection{Week 7}

\textit{Smart contracts --- I didn't see the question in the lecture recording, so I've summarised the required reading below.}
\\
\\
\noindent
A smart contract is an automatically executing program, that is stored on a blockchain.
Users interacting with a smart contract do not need to trust a third party, such as a bank,
because functionality is directly enforced in the code.
When a smart contract is deployed, it is typically immutable and deterministic which ensures an 
authority cannot alter the service that users are running. 
This makes it difficult for continuous development, and bugs can become persistent in a contract.
The article also highlights several potential applications where smart contracts might be valuable, such as:

\begin{itemize}
    \item \textit{NFTs}
    \item \textit{Games}
    \item \textit{Marketplaces}
    \item \textit{Custom digital currencies}
\end{itemize}

\subsection{Week 8}

\textit{Where did you notice the authors assumed a human is still somewhere in the loop?}
\textit{Consider a system where those roles are replaced by AI agents (as explored in ecosystems like NEAR Protocol).}
\textit{Do you think this shift strengthens or weakens the original smart contract thesis? Consider autonomy, verifiability, and security.}
\\
\\
\noindent
In the applications section of the whitepaper,
human interaction in the form of arbiters is assumed in several locations.
For example, human decision about the trustworthiness of different parties is mentioned in multiple places.
In one case, with humans evaluating the trustworthiness of organisations, and in another case determining the reputation of each other.
Additionally, it is mentioned that humans could vote in DAOs to collectively decide how an organisation might allocate its funds.
I think the shift from human input to an AI agent doesn't necessarily strengthen or weaken the original smart contract thesis
but shifts the problem entirely.
There are multiple applications in which an AI arbiter might be a stronger approach, particularly in terms of autonomy.
For example, systems where there may need to be many low impact decisions made in quick succession, 
which might be time consuming for a human but easy for an agent.
Verifiability is another issue for agentic systems, because it is often difficult to determine the validity of the claims of an agent.
This might arise as an issue if one of the parties is not satisfied with the decision of an AI arbiter.

\subsection{Week 9}

\textit{Think of tokens as programmable incentives inside a system.}
\textit{Pick a simple dApp idea (e.g., peer tutoring marketplace, community library, campus carpooling app).}
\\
\\
\noindent
I'll choose the idea of a peer tutoring marketplace.

\begin{itemize}
    \item \textit{What is one specific user action you want to encourage?}
            \\
            \\
            \noindent
            One user action that one would want to encourage is showing up after booking a tutoring session.
            If neither of the users is required to show up after organising a session, it would reduce the utility of the application.

    \item \textit{How does a token cause or enforce that action?}
            \\
            \\
            \noindent
            Using a token in the system provides a monetary incentive for users to provide the
            services that they've advertised. For example, if a tutor gets paid per successful 
            tutoring session, an internal token would incentivise them to show up and tutor.
            Additionally, the student would be incentivised to get the value from their payment
            so they would be more likely to go to a session as well.

    \item \textit{What can users do with the token inside your system (not just trade or sell it)?}
            \\
            \\
            \noindent
            In this system, a student would be able to use the token to purchase services from a tutor.
            A user would be able to advertise a tutoring session in a specific topic at a certain rate.
            This would lead to a marketplace in which tutors could compete to offer the best services in a specific area.
            The competition in this model would hopefully lead to fairer compensation for higher performing teachers.

    \item \textit{If you removed the token, what would break or become significantly worse?}
            \\
            \\
            \noindent
            This system would become significantly worse because it would not allow for the tutors to get paid for 
            their work. This would make it difficult to attract talent, as all the tutors would need to be volunteers.
            This would likely significantly reduce the userbase of the application, both on the student and tutor sides.
            
\end{itemize}
\textit{Keep the answer to approximately one sentence per question.}

\subsection{ Week 10 / Week 11}

\textit{The paper identifies four risk categories:}
\begin{itemize}
    \item \textit{software security}
    \item \textit{transaction costs}
    \item \textit{governance}
    \item \textit{systemic risk}
\end{itemize}
\textit{Which do you think poses the greatest barrier to mainstream DeFi adoption, and why?}
\\
\\
\noindent
This risk category that I believe poses the greatest barrier to mainstream DeFi
adoption is software security.
This is because in smart contracts, bugs are permanent and can often be exploited without defence.
In an economically critical section such as finance, this could cause immeasurable damage.
Governance and transaction cost concerns do raise questions about the relative efficiency of a decentralized approach,
but the software security and continuous development issues are much more serious for potential investors.
The attack surface covers both the exploitation of smart contracts and oracles,
and potential attacks could have a detrimental impact on the entire sector.
Therefore software security issues pose the greatest barrier to broad DeFi adoption.

\end{document}
